% F081: precise mathematical teaching diagram.
\begin{tikzpicture}[x=1cm,y=1cm,>=stealth,font=\small]
\draw[rounded corners=3pt,gray!40] (0,0) rectangle (14.2,9.0);
\begin{scope}[yshift=1cm]
\node at (7.1,7.4) {例2：两个导数零点限定上限，三段异号区间保证存在};
\begin{scope}[shift={(6.1,3.1)},x=2cm,y=0.85cm]
\draw[->,gray!65] (-2.35,0)--(2.4,0) node[right] {$x$};
\draw[->,gray!65] (0,-1.8)--(0,3.9) node[above] {$y$};
\draw[blue!75!black,line width=1.15pt,,domain=-2.1:2.1,samples=160] plot (\x,{(\x)*(\x)*(\x)-3*\x+1});
\fill[blue!75!black] (-2,-1) circle (1.8pt);\node[above left] at (-2,-1) {$(-2,-1)$};
\fill[blue!75!black] (-1,3) circle (1.8pt);\node[above left] at (-1,3) {$(-1,3)$};
\fill[blue!75!black] (1,-1) circle (1.8pt);\node[below right] at (1,-1) {$(1,-1)$};
\fill[blue!75!black] (2,3) circle (1.8pt);\node[above left] at (2,3) {$(2,3)$};
\fill[green!50!black] (-1.8793852415718,0) circle (1.8pt);
\fill[green!50!black] (0.3472963553339,0) circle (1.8pt);
\fill[green!50!black] (1.532088886238,0) circle (1.8pt);
\draw[orange!85!black,line width=1.2pt] (-1.35,3)--(-0.65,3);
\draw[orange!85!black,line width=1.2pt] (0.65,-1)--(1.35,-1);
\end{scope}
\node[] at (7.1,1) {三段异号区间 $(-2,-1),(-1,1),(1,2)$ 各至少一根。};
\end{scope}
\node[text width=13cm,align=center] at (7.1,.9) {$f=x^3-3x+1$，$f^{\prime}$ 只有两个零点，所以原函数恰好三个不同零点。};
\end{tikzpicture}
